\subsection{Available Keys\label{sec:keys}}
\begin{function}{\enverbsetup}
  \begin{syntax}
    \cs{enverbsetup} \marg{key=value}
  \end{syntax}
  This macro can be used to (locally) change the settings of \envlogo\ outside
  of the |\enverb| macro.
\end{function}

The available keys are:
\begin{function}{ignore}
  \begin{syntax}
    ignore = \meta{int}\keydefault{\emph{unused}}
  \end{syntax}
  Set up the number of tokens to ignore at the start of each line. This is only
  applied after the matching |\end| was found, hence doesn't affect the
  |\end|-line. A line containing any non-whitespace tokens but in total
  fewer tokens (counting whitespace) than \meta{int} will lead to undefined
  behaviour.\par
  |ignore| and |auto-ignore| are mutually exclusive.
\end{function}

\begin{function}{auto-ignore}
  \begin{syntax}
    auto-ignore\keydefault{true}
    auto-ignore = \meta{bool}
  \end{syntax}
  If this is true (or no value is given) the number of leading tokens is
  determined by the number of leading tokens found in the line with the matching
  |\end|. This allows for nice nesting of indentation levels in your \LaTeX\
  sources without affecting functionality. A line containing any non-whitespace
  tokens but in total fewer tokens (counting whitespace) than the leading tokens
  in the |\end|-line will lead to undefined behaviour.\par
  |ignore| and |auto-ignore| are mutually exclusive.
\end{function}

\begin{function}{more-ignore}
  \begin{syntax}
    more-ignore = \meta{int}\keydefault{2}
  \end{syntax}
  If |auto-ignore| is true this specifies how many tokens in addition to the
  number of leading tokens shall be ignored at the start of each line. This
  allows for even the contents of the |\enverb|-environment to be further
  indented than the |\begin| and |\end| statement.
\end{function}

\begin{function}{bol}
  \begin{syntax}
    bol = \meta{tokens}\keydefault{\emph{empty}}
    bol+ = \meta{tokens}
    +bol = \meta{tokens}
  \end{syntax}
  Defines which \meta{tokens} are placed at the \textbf{b}egin \textbf{o}f
  \textbf{l}ines. Your \meta{tokens} can be arbitrary contents and are placed as
  is (so can contain macros which work as you intend when using
  \cs{enverbBody}). Keep in mind that if you want to use this with
  \cs{enverbExecute} or \cs{enverbListing} your tokens should be rescanned to
  the same code after they got |\detokenize|d.\par
  The |bol+| variant adds your \meta{tokens} after the current contents of
  |bol|, while |+bol| adds them in front of the contents.
\end{function}

\begin{function}{eol}
  \begin{syntax}
    eol = \meta{tokens}\keydefault{\emph{\string^\string^M\textsubscript{12}}}
    eol+ = \meta{tokens}
    +eol = \meta{tokens}
  \end{syntax}
  Defines which \meta{tokens} are placed at the \textbf{e}nd \textbf{o}f
  \textbf{l}ines. Your \meta{tokens} can be arbitrary contents and are placed as
  is (so can contain macros which work as you intend when using
  \cs{enverbBody}). Keep in mind that if you want to use this with
  \cs{enverbExecute} or \cs{enverbListing} your tokens should be rescanned to
  the same code after they got |\detokenize|d.\par
  The |eol+| variant adds your \meta{tokens} after the current contents of
  |eol|, while |+eol| adds them in front of the contents.
\end{function}

\begin{function}{key-handler}
  \begin{syntax}
    key-handler = \meta{code}\keydefault{\emph{throw an error}}
  \end{syntax}
  If used any undefined keys encountered while parsing the mandatory and
  optional arguments of \cs{enverb} is parsed via \meta{code}. Inside your
  \meta{code} use |#1| to specify where to input the list of unknown keys. This
  can be used to specify additional keys to provide an interface to control your
  wrapping code. This includes the handling of any unknown keys encountered
  while parsing the mandatory \meta{key=value} argument of \cs{enverb}.
\end{function}

\begin{function}{key-set}
  \begin{syntax}
    key-set = \meta{set}
  \end{syntax}
  \envlogo\ uses \pkg{expkv} for its key-parsing. This key is a special form of
  |key-handler| defining a handler that parses the unknown keys as part of
  \meta{set} in \pkg{expkv} (so basically a shorthand for
  |key-handler = \ekvset{|\meta{set}|}{#1}|).
\end{function}

\begin{function}{oarg-not-enverb}
  \begin{syntax}
    oarg-not-enverb\keydefault{false}
    oarg-not-enverb = \meta{bool}
  \end{syntax}
  If this is true (or no value is given) the \meta{key=value} list in the
  optional argument to \cs{enverb} is not parsed as keys of \envlogo\ but only
  inside your defined |key-handler|. Obviously this key doesn't have any effect
  if only used as part of the optional argument of \cs{enverb}.
\end{function}

